\section{Compact path records and protected charts}
\label{sec:protected-charts}

We now record discrete paths without assuming that their rescaled lengths
or parametrizations are tight. Compactness gives subsequential path data.
A protected chart then reconstructs every relevant primal incidence from
finitely many raw word blocks and exact relative lattice heights. Together with the
continuum sewing, the Brownian genericity and raw-word estimates
established above supply the probabilistic inputs to this reconstruction.

\subsection{Compact records of paths}

A path is recorded by its endpoint times, compact range, and cost, together with finitely many ordered marks. This gives a compact state space before an endpoint modulus is available.

Fix a bounded time window \(I=[-T,T]\), a deterministic positive distance denominator \(s_m\), and the grid \(m^{-1}\Z\cap I\). A vertex occurring on a graph path may be represented by any tour time at that vertex in the window. Both endpoint representatives of each used edge are retained. For a path with \(r\) marked consecutive subpaths, record:
\[
 (t_0,\ldots,t_r;\ K_1,\ldots,K_r;\ c_1,\ldots,c_r;\ c).
\]
Here \(t_i\in I\) are the endpoint representatives, \(K_i\) is the finite set of representing times used in the \(i\)-th subpath, including its endpoints, \(c_i\) is an upper bound on its length divided by \(s_m\), and \(c\ge\sum_i c_i\). Costs lie in the compact space \([0,\infty]\). An empty subpath has the singleton range consisting of its endpoint representative. At repeated marks one uses the same chosen endpoint representative. A zero-length change of representative retains both times and their equality as a vertex of the chart graph. Used step identities and equivalences are always those of this graph. Winding and transverse-crossing tests additionally retain their finite ordered marks and open ports; they are not inferred from an unordered compact range alone.

Let \(\mathcal K(I)\) be the space of nonempty compact subsets of \(I\), with Hausdorff distance. Taking the closure of all such records gives a closed subset of
\[
 \mathcal X_{T,r}
 =I^{r+1}\times\mathcal K(I)^r\times[0,\infty]^{r+1}.
\]
An empty set of admissible records is also allowed. We retain these closed sets simultaneously for every integer \(T,r\), every finite list of rational time restrictions on the marks and ranges, and every finite chart graph defined below.

\begin{lemma}[Certificate compactness]\label{bg:certificates}
These observations take values in a compact metrizable product space. Along any sequence for which the rescaled contour laws converge, they admit a joint subsequential limit with the contours. In a joint almost surely convergent realization:
\begin{enumerate}[label=(\roman*)]
\item Every limiting finite certificate is approximated by actual discrete path records obeying any prescribed open loss in its endpoint, range, and cost restrictions.
\item Every sequence of admissible discrete records has a subsequential limiting certificate.
\item Finitely many ordered subpaths of one discrete path may be extracted simultaneously. If its limiting total cost is finite, the sum of the limiting subpath costs is no larger than that total cost.
\item The projected range of a limiting path is connected. Its endpoints are the projections of its recorded endpoint times.
\end{enumerate}
There is no concatenation assertion for independently extracted certificates.
\end{lemma}
\begin{proof}
The hyperspace of closed subsets of a compact metric space is compact, with the empty set adjoined as an isolated point. The countable product is compact and metrizable. Joint tightness with the contours and a Skorokhod realization therefore give convergence of all the closed observation sets.

For a point of a limiting closed set, Hausdorff convergence supplies points in the approximating closed sets. A second approximation replaces those points by actual records. Any open restriction containing the limiting record is eventually satisfied. This proves (i). Compactness proves (ii). Recording the marked subpaths as one tuple proves (iii); addition is continuous where the finite total cost bounds all entries, and the inequalities pass to the limit there.

For (iv), consecutive path vertices have asymptotically equal projected locations, by the contour chord rules and continuity of \(\eta\); the same holds for different representatives of one vertex. Join their projected locations inside balls of vanishing radius. The resulting connected compact sets have Hausdorff limits equal to the projected recorded ranges. Hausdorff limits of nonempty connected compact sets are connected. This argument is purely topological and puts no bound on the length of the small joining pieces.
\end{proof}

A spatial confinement condition always means compact confinement in an open set, or confinement after an arbitrarily small open enlargement. This convention is needed because limiting records are closed data. For example, a crossing from an inner circle to an outer circle is recorded with a margin on both sides. First and last passages through intermediate bands are then available from the simultaneous subpath records. A limiting certificate is \emph{attainable} in this sense of approximation with open loss.

The eventual internal passage rule \(F_U(x,y)\) is the infimum of finite costs of attainable certificates with endpoints \(x,y\) and range compactly contained in \(U\). At this stage it is only a passage functional. Its concatenation, continuity, and comparison with \(D_h\) are proved later from small circuits and open ports.

\begin{definition}[Universal endpoint observations]\label{bg:universal-tests}
For each finite chart graph, retain in addition the following countable family of variables. Given two finite unions \(J_1,J_2\) of rational time intervals, select vertices having at least one representative in the respective unions. A rational chart confinement specifies the allowed vertex and step indices by a finite union of such intervals; the confined graph consists of the specified vertices and those specified edges whose endpoints are both allowed. Take the supremum over the selected endpoint pairs of their minimal distance in this confined graph, divided by \(s_m\). A disconnected pair has cost infinity; an empty supremum is zero. Also retain finite refinements with several endpoint sets and several confinements.
\end{definition}

These variables lie in a countable product of \([0,\infty]\) and hence preserve compactness. They are measurable observations of the finite chart graph. They are retained independently of existential path certificates: convergence of the latter alone does not transfer a statement uniform over arbitrary approaching discrete endpoints. Once a time or spatial endpoint modulus is proved, rational time covers and open confinement losses let it be expressed using \Cref{bg:universal-tests}.

\subsection{Words, instructions, and stable comparisons}

Before specifying a chart, we show how finitely many external flexible
choices can be fixed from its limiting contour. In this subsection the
lattice time unit is $N\to\infty$: times are divided by $N$ and heights
by $\sqrt N$.

Write the inventory alphabet as
$\mathcal A=\{h,c,H,C,F\}$, where lower-case letters produce a burger,
upper-case $H,C$ order the corresponding type, and $F$ orders the freshest
available burger.  Burger identities include their creation times.  A
\emph{local reduction} of a finite word starts with two empty stacks,
performs every possible match with a burger produced in the word, and
records an order which cannot be filled locally.  An externally unresolved
$F$ is assigned a type only after this reduction.  A finite instruction
$\sigma$ specifies that type as a function of a finite rational partition
of time.  Denote the resulting two-coordinate increment path by
$Z_N^{\sigma}$.  Its net increment is an exact integer function of the
raw word and $\sigma$.

\begin{lemma}[Input-independent local reduction]\label{loc:local-reduction}
The matches made by local reduction are independent of the stacks supplied
before the word.  Whenever an $F$ is externally unresolved, every available
burger at that step was produced before the beginning of the word.
Consequently its two actual candidates, if both exist, precede that
beginning.
\end{lemma}
\begin{proof}
A burger produced inside the word is newer than every input burger.
Whenever a compatible locally produced burger is available, a rigid order
uses the newest such burger of its type, and an $F$ uses the newest of all
locally produced burgers.  These choices do not depend on the input.
Induction through the word proves both assertions.  In particular, removing
an input burger at an externally unresolved order does not change either
local stack.
\end{proof}

The Brownian genericity proved earlier gives the stable comparison needed
for external flexible decisions.  A coordinate-$i$ horizontal pair is a pair $s<t$ such that
$z_i(s)=z_i(t)=\inf_{[s,t]}z_i$.  A triple means three distinct times with
the same coordinate value and no lower value between them.  For the
correlated Brownian motion in this paper, and for its restrictions under
interior excursion absolute continuity, almost surely:
\begin{enumerate}[label=(\roman*)]
\item there is no nontrivial horizontal pair common to both coordinates;
\item a time belonging to a coordinate triple is not the endpoint of an
additional nontrivial horizontal pair in the other coordinate.
\end{enumerate}
Lemma~\ref{arm:generic-chords} proves both assertions from Brownian wedge
survival and independent-gap density. Its hypothesis is $d_2>1/2$; no
bound on $d_\gamma$ is needed.

\begin{lemma}[Uniform finite instructions]\label{loc:finite-instructions}
Let $z=(z_1,z_2)$ be a continuous path on $[A,v]$ satisfying (i)--(ii)
above.  Fix $A<u-\delta<u\le v$, with $\delta>0$.  For $\epsilon\ge0$,
let $K_\epsilon$ consist of triples
$(s_1,s_2,t)\in[A,u-\delta]^2\times[u,v]$ such that, for $i=1,2$,
\[
 |z_i(s_i)-z_i(t)|\le\epsilon,
 \qquad
 \inf_{[s_i,t]}z_i\ge z_i(t)-\epsilon.
\]
There are $\epsilon_0>0$, a finite rational partition of $[u,v]$, and a
type assigned to each partition interval such that the assignment equals
the sign of $s_1-s_2$ for every triple in $K_{\epsilon_0}$.  Partition
boundaries may be placed in open time gaps containing no such triple.
The same instruction is valid for every path sufficiently close to $z$
in the uniform norm, with a sufficiently small additional lattice height
error.
\end{lemma}
\begin{proof}
The sets $K_\epsilon$ are compact and decrease to $K_0$ as
$\epsilon\downarrow0$.  Property (i) implies $s_1\ne s_2$ on $K_0$,
so $|s_1-s_2|$ has a positive minimum there if $K_0$ is nonempty.
Let $T_+$ and $T_-$ be the projections to $t$ of the two closed sign
parts of $K_0$.  They are compact.  If a time belonged to both, the
supplier in at least one coordinate would have two distinct choices.
Those choices together with $t$ would make a scalar triple, while the
other coordinate has a horizontal pair ending at $t$, contrary to (ii).
Thus $T_+$ and $T_-$ are disjoint compact subsets of the time interval.
Choose disjoint finite unions of open intervals with rational endpoints
around them.  Decreasing $\epsilon_0$ if necessary makes every triple
in $K_{\epsilon_0}$ belong to the corresponding sign neighborhood;
otherwise a sequence of contrary triples has a limit in $K_0$.
The resulting finite partition gives the assertion.  Uniform perturbation
of the path only increases the two errors in the definition of
$K_\epsilon$ by a fixed multiple of that perturbation.
\end{proof}

\begin{lemma}[Correctness of surrogate resolutions]\label{loc:first-error}
Consider finitely many ordinary word intervals, each with a left extension
of positive limiting duration.  Suppose all suppliers of externally
unresolved orders in the used portions lie in a fixed bounded time window,
and that their distance from the used time is at least a fixed positive
number in rescaled units.  Fix a generic limiting contour and instructions
from Lemma~\ref{loc:finite-instructions}.  For all sufficiently large
$N$, every surrogate resolved word in a sufficiently small contour tube
about that limit has the actual freshest-burger resolutions on the used
portions, provided its initial stacks are those represented by its past
contour, its unused left part is processed with actual or already
verified context-imposed resolutions, and no empty-stack order is imposed.
\end{lemma}
\begin{proof}
If a first incorrect decision exists, all preceding resolutions are
correct.  Its actual candidates can therefore be read from the already
correct past contour.  By Lemma~\ref{loc:local-reduction}, both candidates
precede the extension.  They form approximate horizontal pairs with height
error at most one lattice unit, hence belong to $K_{\epsilon_0}$ after
rescaling.  Lemma~\ref{loc:finite-instructions} gives exactly the
freshness comparison prescribed by the instruction, a contradiction.
Apply this argument successively to the finitely many portions.  Supplier
truncation is robust: a left cut whose two contour heights are strictly
below all tested heights forces every required backward first hit to occur
after the cut.  Strict positivity supplies the no-empty-stack condition
in a finite interior excursion chart.
\end{proof}

In this induction the preceding contour is the actual contour, or the
contour of a globally imposed resolution whose preceding decisions have
already been verified.  A locally chosen instruction for a used portion
does not assert correctness for externally unresolved orders confined to
its unused extension.  The distinction will matter in anchoring the
reference heights in Section~\ref{loc:section}.

\subsection{Protected finite charts}

A chart must preserve exact incidences, even where contour convergence alone does not decide a lattice equality. We specify its raw blocks, relative heights, and finite instructions first; \Cref{sec:local-kernels} will then determine the conditional law of its passage observations.

Choose nested independent Poisson time marks, and round discrete cuts to even integers. For a spatial compact set \(K\) inside an open buffer \(U\), retain finitely many marked \emph{used} intervals whose curve ranges lie in \(U\), together with unused left extensions when needed. Adjacent intervals may be merged. All preimages of \(K\), and of the specified compact neighborhood of \(K\), must be covered by interiors of used intervals. Such finite covers exist: transience bounds their time window, the preimage of a compact set is closed in that window, and continuity allows a sufficiently fine Poisson mesh to cover it by intervals lying in \(U\). All exact forest heights below are anchored at used starts. The raw extension is retained to compute internal matches, but its externally unmatched flexible orders may have fixed surrogate types; its genuine resolved increment is never silently substituted for that surrogate in an exact height identity.

For each scalar contour, a cut presents at most two boundary lineages. Their heights are measured relative to the start of their retained block. In addition to the raw words and cut lengths, the chart records:
\begin{enumerate}[label=(\alph*)]
\item the relative start-height differences needed to compare retained lineages, encoded by a forest in each coordinate;
\item finite unions of allowed height bands for each pair of boundary sides, and any minima of retained intervening blocks that can obstruct a candidate match; every band endpoint is either a strict rational threshold or a minimum computed from a specified retained word;
\item directed incoming/outgoing side labels and the finite relative order of the local traversals;
\item a finite time partition with the decisions for flexible orders not resolved within the extended raw block.
\end{enumerate}
An unused absolute height completing the forest coordinates has no geometric role. Feasible relative differences lie on the projection of the parity lattice: at even times the coordinate sum is even. Heights inside an individual block are then computed from its raw word and its prescribed external decisions.

In an intrinsic chart, each time is labeled by its block and its time relative to that block's start. The compact record spaces above are used with the finite disjoint union of these closed block intervals in place of \(I\); the same hyperspace argument applies. Rational endpoint and confinement tests use these block coordinates. Absolute starts on the global tour and omitted-gap durations are exterior data. Conditional on the full contour they push intrinsic records into the global time window, but they are not inputs to a spatially local kernel.

For a fixed finite choice of these instructions and lattice parameters, the resulting graph and all its certificate and universal endpoint observations are functions of the retained raw words. The instructions are part of the input to the reference experiment; their validity in an actual word is proved below and in the guard construction.

For a scalar coordinate \(X\) and a retained block \(I_i=[a_i,b_i]\), put \(m_i=\min_{I_i}X\). Its left and right \emph{ports} are the contour-tree lineages from the visits at \(a_i\) and \(b_i\) down to height \(m_i\). They are compact intervals of length coordinates; their spatial images are contained in \(\eta(I_i)\). A contact with an earlier block lies on the left port, and one with a later block lies on the right port. Several times representing the same port height are distinguished inside the retained raw block.

Local field data recover scalar contacts through the two labelled families
of frontier arcs. Retaining the labels is essential: an equality of spatial
points can also be caused by the other contour coordinate.

\begin{lemma}[Local coloured port data]\label{bg:local-ports}
For finitely many retained blocks whose closures lie in a buffered local traversal, the local fields and the traversal's side labels determine the following data: the scalar port relations in the protected region, the connected components of the port-overlap graph, and the relative block-start heights within each such component. The only order information required in this statement is the finite relative order of the directed local traversals.
\end{lemma}
\begin{proof}
Use the labelled frontier strands in \Cref{bg:continuum-inputs}(iv). In the imaginary-geometry construction, strands of the same labelled family merge on meeting. Their shared arcs give scalar contour gluing with the corresponding side convention; see \cite[Section~2.1.3 and Lemma~2.4]{GwynneMillerSheffieldHarmonic}. The block boundaries are the first and last frontier portions of the clocked local traversal, so these labels do not require an absolute exterior contour height.

For each pair of ports, take the closure of their correspondence on nontrivial common arcs of the \emph{same} labelled family, with the incoming and outgoing sides specified. Initially this is a relation on their length-coordinate intervals. For an ordered block pair \(i<j\), pull it back to endpoint times by the right-port map \(s\mapsto X(b_i)-X(s)\) and the left-port map \(t\mapsto X(a_j)-X(t)\); restrict to times passing the retained suffix and prefix minimum tests, which ensure membership in these ports. Denote this pulled-back closed relation by \(E_{ij}\). Every such common arc gives a scalar contact. The scalar equal-minimum relation is closed, so taking the closure adds no false scalar contacts. Conversely, a scalar contact between interior times of two retained blocks is approached by ordinary scalar chords with endpoints in those same blocks. For an ordinary excursion pair use nearby higher levels. At a branch triple use higher levels for the adjacent pairs and lower levels for the outer pair. Choose levels avoiding the countable set of scalar local-minimum heights. Strict Brownian local minima and their distinct heights justify these approximations. Thus every protected scalar contact belongs to \(E_{ij}\). This uses \Cref{arm:generic-chords} for the Brownian sewing event and the scalar approximation argument detailed in \Cref{bg:chord-approximation}.

The relation is not defined by arbitrary equality of spatial points: that equality can be caused through the other colour. Keeping the labelled common arcs is what distinguishes the two relations. Two tree lineages intersect in a lineage, so the common length coordinates of one pair of ports form an interval. There are finitely many ports. Consequently their relations have finitely many such pieces, even if the associated surfaces have infinitely many beads or the strands have infinitely many individual points of contact.

The real field gives the quantum-length coordinates on the retained arcs. Evaluating the two coordinates at a common arc point determines a single additive difference of block-start heights. A spanning forest of the nontrivial-overlap graph stores all differences in each connected component. A protected isolated endpoint contact also joins the same component, by the ordinary-chord approximation just proved. No difference between distinct components is recovered or needed.

These are measurable local data. The strands, their side labels, local clock, and length coordinates are measurable by the cited construction. Their common-arc correspondence is a measurable closed relation: it can be tested on compact parameter subintervals, with positive common length exhausted by rational lower bounds. Standard measurable selection on these compact relations chooses a common point; every choice gives the same additive difference. Under the countable family of reference laws a common completed-measurable version has a Borel version, by using their probability-weighted sum. Choosing a spanning forest by the fixed finite order is then measurable.
\end{proof}

The relative heights have now been recovered within each overlap component.
A possible contact can still be blocked by a minimum in an omitted time
gap. The following clearance argument replaces that exterior test by a
rational threshold chosen from the local port relation.

\begin{lemma}[Clearance from untested minima]\label{bg:omitted-clearance}
Let \(K_*\) be a compact spatial set whose full preimage lies in the interiors of finitely many retained blocks in a bounded time window. Fix a coordinate and two blocks \(i<j\) in the same overlap component \(\mathcal A\). Translate its component-root height to zero. Test exactly the suffix of block \(i\), the prefix of block \(j\), and every intervening retained block in that component. The combined untested set is
\[
 G_{ij}=\bigcup_{\ell=i}^{j-1}[b_\ell,a_{\ell+1}]
 \ \cup\!\!\bigcup_{\substack{i<\ell<j\\\ell\notin\mathcal A}} I_\ell
 \ \subset[b_i,a_j].
\]
Thus it contains precisely the intervening omitted gaps, their cut endpoints, and the intervening blocks in other components. Set \(m_{ij}=\min_{G_{ij}}X\).

Let \(C_{ij}\) consist of pairs \((s,t)\) in the two blocks with \(\eta(s),\eta(t)\in K_*\), equal endpoint height, and every tested height at least that endpoint height. Its height set \(H_{ij}\) is compact and
\[
 \operatorname{dist}(H_{ij},\{m_{ij}\})>0
\]
whenever \(C_{ij}\ne\varnothing\). There is a rational band test, selected using only the local port data and tested retained pieces, which decides every protected candidate's untested minimum condition. It is correct uniformly for sufficiently small contour perturbations and bounded multiples of a vanishing height-mesh error.
\end{lemma}
\begin{proof}
Interval minima are continuous in their endpoints, so \(C_{ij}\) and \(H_{ij}\) are compact. Suppose a candidate has height \(m_{ij}\), and let \(\tau\in G_{ij}\) minimize the untested contour. All heights between the candidate endpoints are then at least \(m_{ij}\). Hence \(s,t,\tau\) belong to one scalar class and \(\eta(\tau)=\eta(s)\in K_*\). An omitted gap or cut cannot contain such a visit, by full-preimage retention. If \(\tau\) belongs to another retained block, its scalar contact with \(s\), approximated by common arcs as in \Cref{bg:local-ports}, puts that block in the same overlap component. This too contradicts its inclusion in \(G_{ij}\). Thus \(m_{ij}\notin H_{ij}\), proving positive separation.

The untested minimum is used to prove separation, but is not an input to the local selection. Instead intersect \(C_{ij}\) with the local relation \(E_{ij}\) from \Cref{bg:local-ports}. Call the resulting candidates good, and the other candidates bad. By the preceding separation, their respective height sets satisfy
\[
 H_{\rm good}\subset(-\infty,m_{ij}-\delta],\qquad
 H_{\rm bad}\subset[m_{ij}+\delta,\infty)
\]
for some \(\delta>0\). Both sets are compact, including the bad set. If both are nonempty, choose the first rational \(\alpha\) in a fixed enumeration with
\(\sup H_{\rm good}<\alpha<\inf H_{\rm bad}\), and allow heights below \(\alpha\). If only bad candidates exist, allow none; if only good candidates exist, choose \(\alpha>\sup H_{\rm good}\). An empty candidate set requires no test. All these choices use only local data and the finite order. Extrema and the first-admissible-rational choice are measurable; equivalently exhaust complements of closed relations by positive rational distance and take compact-set extrema.

For stability, suppose uniformly convergent contours, convergent cuts and a vanishing height mesh gave incorrectly classified protected lattice candidates. After taking endpoint subsequences, all exact retained tests pass to the limiting weak tests, so the limit lies in \(C_{ij}\). The combined untested minima converge to \(m_{ij}\). At every such limiting candidate both its relation to \(m_{ij}\) and its classification by the chosen band have strict and agreeing signs. This contradicts incorrect classification. The same argument covers the empty-good or empty-bad case. It proves correctness throughout one sufficiently small contour tube, rather than only along the originally sampled sequence.
\end{proof}

We can now combine the port relations, their strict clearance, and the
finite flexible-order instructions into an exact graph reconstruction.

\begin{lemma}[Protected chart reconstruction]\label{bg:protected-atlas}
On the Brownian sewing event specified in \Cref{arm:generic-chords}, compactly confined graph observations can be exhausted by the finite charts just described. More precisely, for \(K\Subset U\), there is a countable collection of buffered chart descriptions with the following properties.
\begin{enumerate}[label=(\roman*)]
\item Every graph path with projected range in \(K\), and every finite collection of its marked subpaths, is eventually represented in one of these charts. The same holds for the universal endpoint observations with an open confinement loss.
\item Within the compact confinement, the reconstructed chart has exactly the primal incidences of the original graph for all sufficiently large scales. No unrecorded minimum in an omitted gap is used to decide one of these incidences.
\item The continuous block input, relevant forest differences, and allowed local gluing bands are measurable from the directed local traversals, frontier lengths, and the finite global traversal order. Exterior flexible decisions are treated by the buffered comparison instructions of \Cref{loc:finite-instructions,loc:first-error}.
\item Disjoint spatial buffers use disjoint raw blocks and distinct local observations. The countable atlas can be refined to accommodate any prescribed finite collection of observations.
\end{enumerate}
\end{lemma}
\begin{proof}
Choose nested neighborhoods
\[
 K\subset\operatorname{int}K_*\Subset U_0\Subset U_1\Subset U.
\]
Cover the full preimage of \(K_*\) by used intervals whose entire ranges have closure in \(U_0\), and give them left extensions whose entire ranges have closure in \(U_1\). Transience supplies a bounded enclosing window. For a bilateral discrete application enlarge the window first to coordinate descent cuts strictly below every tested height, so all needed partners are in the window. This removes the possibility of an unrecorded supplier escaping to infinity. Record every directed \(U_1\)-traversal reaching \(\overline U_0\). There are finitely many: each crosses a fixed positive distance, and uniform continuity on the bounded time window bounds its duration below. The dense local-traversal construction in the proof of \cite[Lemma~2.4]{GwynneMillerSheffieldHarmonic} includes them all. The real field needed to clock them is restricted to \(U_1\subset U\); a slightly larger imaginary buffer suffices to reconstruct them.

Apply \Cref{bg:local-ports} separately in the two coordinates, storing each component's offsets in a forest. Apply \Cref{bg:omitted-clearance} to each pair of blocks in a component. Test intervening blocks only within that component; all other blocks belong to the combined untested set in that Lemma. No protected incidence can join different components for infinitely many scales, since its limiting scalar contact would join them by \Cref{bg:local-ports}. There are finitely many block pairs, so one contour tube works for all their band tests.

Keep the exact discrete sewing dictionary at the retained pieces. Primal contour-vertex identifications use the appropriate equal-height weak-minimum tests. An edge supplied by the other coordinate uses a genuinely matched up-step and down-step, with the strict one-step-above-base condition in its excursion interior and the prescribed endpoint corner offsets. These tests are functions of the actual raw steps and exact lattice forest shifts. In particular a weak horizontal pair in the other coordinate is not by itself declared to be a primal edge. For example, the scalar word with heights \(0,1,0,1,0\) has a weak pair between its first and last visits but two separate matched excursions. Every retained suffix, prefix and intervening minimum, including a critical branch minimum, is evaluated exactly in the discrete word. Only an \emph{untested} obstruction is replaced by a rational band. \Cref{bg:omitted-clearance} permits the bounded height-mesh errors from the step convention and proves that this replacement preserves each relevant incidence.

Exterior flexible resolutions require another step. A flexible order unresolved within its extended raw block, but lying in its used portion, has both candidate suppliers before the extension and hence at a positive limiting time separation from the used portion. Each positive-span supplier chord has the same spatial image as its decision. Since the whole used image lies in \(U_0\), every supplier's local traversal belongs to the finite \(U_1\)-traversal family just retained. The no-common-chord and mixed-branch exclusions of \Cref{arm:generic-chords} separate the two possible supplier orders into compact disjoint sets of decision times. The finite rational instructions and first-error proof in \Cref{loc:finite-instructions,loc:first-error} therefore make the resolved \emph{used} steps exact on a sufficiently small tube. No recursively enlarged collection of raw supplier blocks is needed: their continuum traversal labels and finite order supply the comparison. Externally unmatched flexible orders in the unused extension need not be genuinely typed to compute the internal matches or used increments. Their surrogate paths enter only the macroscopic reference input, using \Cref{arm:uniform-past-limit}; the exact density identities of \Cref{sec:local-kernels} keep genuine and surrogate extension increments distinct. Anchoring exact heights at used starts prevents these unused types from changing any retained incidence.

We have reconstructed every primal incidence used by a path or an ordered subpath record confined to \(K\), with a margin inside \(K_*\). Rational intrinsic time covers with image in \(K_*\) give the corresponding confined chart graphs for the universal endpoint tests. Their spatial interpretation uses the open loss specified above. Poisson refinements, rational bands, finite instruction lists and finite traversal orders give countably many descriptions; their durations, paths and forest shifts remain continuous or lattice input variables. Refinements accommodate any finite collection of observations. Disjoint spatial buffers have disjoint time preimages, and the unused extensions can be trimmed inside the respective buffers. Thus their retained raw blocks are disjoint, proving all four assertions.
\end{proof}

The passage from this deterministic chart description to conditional local kernels is not automatic. In particular, relative lattice shifts cannot be replaced by continuous offsets solely because the contour converges. \Cref{sec:local-kernels} supplies a lattice local limit, guards external flexible decisions, and identifies the conditional kernels by a reference-block change of density.
